
\documentclass{article}
\usepackage{colortbl}
\usepackage{makecell}
\usepackage{multirow}
\usepackage{supertabular}

\begin{document}

\newcounter{utterance}

\centering \large Interaction Transcript for game `hot\_air\_balloon', experiment `air\_balloon\_survival\_en\_complexity\_hard', episode 1 with kimi{-}k2.5{-}without{-}reasoning{-}t1.0.
\vspace{24pt}

{ \footnotesize  \setcounter{utterance}{1}
\setlength{\tabcolsep}{0pt}
\begin{supertabular}{c@{$\;$}|p{.15\linewidth}@{}p{.15\linewidth}p{.15\linewidth}p{.15\linewidth}p{.15\linewidth}p{.15\linewidth}}
    \# & $\;$A & \multicolumn{4}{c}{Game Master} & $\;\:$B\\
    \hline

    \theutterance \stepcounter{utterance}  
    & & \multicolumn{4}{p{0.6\linewidth}}{
        \cellcolor[rgb]{0.9,0.9,0.9}{
            \makecell[{{p{\linewidth}}}]{
                \texttt{\tiny{[P1$\langle$GM]}}
                \texttt{You are participating in a collaborative negotiation game.} \\
\\ 
\texttt{Together with another participant, you must agree on a single set of items that will be kept. Each of you has your own view of how much each item matters to you (importance). You do not know how the other participant values the items. Additionally, you are given the effort each item demands.} \\
\texttt{You may only agree on a set if the total effort of the selected items does not exceed a shared limit:} \\
\\ 
\texttt{LIMIT = 8145} \\
\\ 
\texttt{Here are the individual item effort values:} \\
\\ 
\texttt{Item effort = \{"A08": 963, "B95": 204, "A39": 418, "A65": 961, "B81": 238, "B34": 623, "A28": 44, "C82": 633, "B69": 247, "A87": 645, "B64": 230, "B15": 250, "C79": 731, "B16": 405, "C53": 389, "C83": 216, "C10": 637, "A97": 156, "A25": 738, "C84": 307, "A71": 762, "B23": 737, "B96": 901, "C51": 369, "C17": 2, "B19": 730, "C72": 722, "A94": 704, "A46": 315, "B43": 455, "B37": 510, "A88": 175, "A96": 692, "C60": 150, "A75": 32\}} \\
\\ 
\texttt{Here is your personal view on the importance of each item:} \\
\\ 
\texttt{Item importance values = \{"A08": 1059, "B95": 300, "A39": 514, "A65": 1057, "B81": 334, "B34": 719, "A28": 140, "C82": 729, "B69": 343, "A87": 741, "B64": 326, "B15": 346, "C79": 827, "B16": 501, "C53": 485, "C83": 312, "C10": 733, "A97": 252, "A25": 834, "C84": 403, "A71": 858, "B23": 833, "B96": 997, "C51": 465, "C17": 98, "B19": 826, "C72": 818, "A94": 800, "A46": 411, "B43": 551, "B37": 606, "A88": 271, "A96": 788, "C60": 246, "A75": 128\}} \\
\\ 
\texttt{Goal:} \\
\\ 
\texttt{Your goal is to negotiate a shared set of items that benefits you as much as possible (i.e., maximizes total importance to YOU), while staying within the LIMIT. You are not required to make a PROPOSAL in every message {-} you can simply negotiate as well. All tactics are allowed!} \\
\\ 
\texttt{Interaction Protocol:} \\
\\ 
\texttt{You may only use the following structured formats in a message:} \\
\\ 
\texttt{PROPOSAL: \{'A', 'B', 'C', …\}} \\
\texttt{Propose keeping exactly those items.} \\
\texttt{REFUSE: \{'A', 'B', 'C', …\}} \\
\texttt{Explicitly reject opponent's proposal.} \\
\texttt{ARGUMENT: \{'...'\}} \\
\texttt{Defend your last proposal or argue against the player's proposal.} \\
\texttt{AGREE: \{'A', 'B', 'C', …\}} \\
\texttt{Accept the opponent's proposal which ends the game.} \\
\texttt{STRATEGIC REASONING: \{'...'\}} \\
\texttt{	Describe your strategic reasoning or anticipation explaining your choice of action. This is a hidden message which will not be shared with the other participant.} \\
\\ 
\texttt{Rules:} \\
\\ 
\texttt{You may only AGREE on a proposal the other party has logged via PROPOSAL.} \\
\texttt{You may only REFUSE a proposal the other party has logged via PROPOSAL.} \\
\texttt{Total effort of any PROPOSAL or AGREE set must be ≤ LIMIT.} \\
\texttt{Do NOT reveal your hidden importance scores.} \\
\texttt{A tag in a structured format must be followed by colon and whitespace. The argument must be a python set containing 0 or more strings.} \\
\texttt{So, it must be of the form TAG: \{...\}} \\
\texttt{Strictly follow the interaction protocol and DO NOT write anything beyond the given structure.} \\
\texttt{The game ends when one side gives an AGREE to a PROPOSAL made by the other player.} \\
\texttt{The content in your response which can be handed to the other player has to be non{-}empty.} \\
\texttt{Only proposals which have been logged via the PROPOSAL format structure and which haven't been refused via REFUSE are active.} \\
\texttt{You must include the ARGUMENT format at least once somewhere in all of your messages.} \\
\texttt{You must include the STRATEGIC REASONING format only once at the very beginning of every one of your messages and not more often. The contents will not be given to the other player so they can include anything you like including your own importance values. Here you should reason step by step to come up with you next move.} \\
            }
        }
    }
    & \\ \\

    \theutterance \stepcounter{utterance}  
    & \multicolumn{4}{p{0.6\linewidth}}{
        \cellcolor[rgb]{0.9,0.9,0.9}{
            \makecell[{{p{\linewidth}}}]{
                \texttt{\tiny{[P1$\rangle$GM]}}
                \texttt{STRATEGIC REASONING: \{'I need to maximize my total importance while staying under the 8145 effort limit. Let me analyze my items by importance/effort ratio to find efficient selections. High value items for me include: A08 (1059/963=1.10), A65 (1057/961=1.10), B96 (997/901=1.11), A71 (858/762=1.13), A25 (834/738=1.13), B23 (833/737=1.13), B19 (826/730=1.13), C72 (818/722=1.13), A94 (800/704=1.14), A96 (788/692=1.14), C79 (827/731=1.13). I should aim for a proposal that includes many of these high{-}value items. Let me calculate a strong opening proposal: A08+A65+B96+A71+A25+B23+B19+C72+A94+A96+C79 = 963+961+901+762+738+737+730+722+704+692+731 = 7941 effort, importance = 1059+1057+997+858+834+833+826+818+800+788+827 = 8697. That leaves 204 room. I could add B95 (204 effort, 300 importance) to get 8145 exactly! Total importance would be 8997. This seems like an excellent opening proposal.'\}} \\
\\ 
\texttt{PROPOSAL: \{'A08', 'A65', 'B96', 'A71', 'A25', 'B23', 'B19', 'C72', 'A94', 'A96', 'C79', 'B95'\}} \\
\\ 
\texttt{ARGUMENT: \{'This proposal uses exactly 8145 effort, the maximum allowed, ensuring we get full value from our shared limit. It includes 12 items that represent a balanced selection across different categories.'\}} \\
            }
        }
    }
    & & \\ \\

    \theutterance \stepcounter{utterance}  
    & & \multicolumn{4}{p{0.6\linewidth}}{
        \cellcolor[rgb]{0.9,0.9,0.9}{
            \makecell[{{p{\linewidth}}}]{
                \texttt{\tiny{[GM$\rangle$P2]}}
                \texttt{You are participating in a collaborative negotiation game.} \\
\\ 
\texttt{Together with another participant, you must agree on a single set of items that will be kept. Each of you has your own view of how much each item matters to you (importance). You do not know how the other participant values the items. Additionally, you are given the effort each item demands.} \\
\texttt{You may only agree on a set if the total effort of the selected items does not exceed a shared limit:} \\
\\ 
\texttt{LIMIT = 8145} \\
\\ 
\texttt{Here are the individual item effort values:} \\
\\ 
\texttt{Item effort = \{"A08": 963, "B95": 204, "A39": 418, "A65": 961, "B81": 238, "B34": 623, "A28": 44, "C82": 633, "B69": 247, "A87": 645, "B64": 230, "B15": 250, "C79": 731, "B16": 405, "C53": 389, "C83": 216, "C10": 637, "A97": 156, "A25": 738, "C84": 307, "A71": 762, "B23": 737, "B96": 901, "C51": 369, "C17": 2, "B19": 730, "C72": 722, "A94": 704, "A46": 315, "B43": 455, "B37": 510, "A88": 175, "A96": 692, "C60": 150, "A75": 32\}} \\
\\ 
\texttt{Here is your personal view on the importance of each item:} \\
\\ 
\texttt{Item importance values = \{"A08": 1059, "B95": 300, "A39": 514, "A65": 1057, "B81": 334, "B34": 719, "A28": 140, "C82": 729, "B69": 343, "A87": 741, "B64": 326, "B15": 346, "C79": 827, "B16": 501, "C53": 485, "C83": 312, "C10": 733, "A97": 252, "A25": 834, "C84": 403, "A71": 858, "B23": 833, "B96": 997, "C51": 465, "C17": 98, "B19": 826, "C72": 818, "A94": 800, "A46": 411, "B43": 551, "B37": 606, "A88": 271, "A96": 788, "C60": 246, "A75": 128\}} \\
\\ 
\texttt{Goal:} \\
\\ 
\texttt{Your goal is to negotiate a shared set of items that benefits you as much as possible (i.e., maximizes total importance to YOU), while staying within the LIMIT. You are not required to make a PROPOSAL in every message {-} you can simply negotiate as well. All tactics are allowed!} \\
\\ 
\texttt{Interaction Protocol:} \\
\\ 
\texttt{You may only use the following structured formats in a message:} \\
\\ 
\texttt{PROPOSAL: \{'A', 'B', 'C', …\}} \\
\texttt{Propose keeping exactly those items.} \\
\texttt{REFUSE: \{'A', 'B', 'C', …\}} \\
\texttt{Explicitly reject opponent's proposal.} \\
\texttt{ARGUMENT: \{'...'\}} \\
\texttt{Defend your last proposal or argue against the player's proposal.} \\
\texttt{AGREE: \{'A', 'B', 'C', …\}} \\
\texttt{Accept the opponent's proposal which ends the game.} \\
\texttt{STRATEGIC REASONING: \{'...'\}} \\
\texttt{	Describe your strategic reasoning or anticipation explaining your choice of action. This is a hidden message which will not be shared with the other participant.} \\
\\ 
\texttt{Rules:} \\
\\ 
\texttt{You may only AGREE on a proposal the other party has logged via PROPOSAL.} \\
\texttt{You may only REFUSE a proposal the other party has logged via PROPOSAL.} \\
\texttt{Total effort of any PROPOSAL or AGREE set must be ≤ LIMIT.} \\
\texttt{Do NOT reveal your hidden importance scores.} \\
\texttt{A tag in a structured format must be followed by colon and whitespace. The argument must be a python set containing 0 or more strings.} \\
\texttt{So, it must be of the form TAG: \{...\}} \\
\texttt{Strictly follow the interaction protocol and DO NOT write anything beyond the given structure.} \\
\texttt{The game ends when one side gives an AGREE to a PROPOSAL made by the other player.} \\
\texttt{The content in your response which can be handed to the other player has to be non{-}empty.} \\
\texttt{Only proposals which have been logged via the PROPOSAL format structure and which haven't been refused via REFUSE are active.} \\
\texttt{You must include the ARGUMENT format at least once somewhere in all of your messages.} \\
\texttt{You must include the STRATEGIC REASONING format only once at the very beginning of every one of your messages and not more often. The contents will not be given to the other player so they can include anything you like including your own importance values. Here you should reason step by step to come up with you next move.} \\
\\ 
\texttt{You will now receive the first message of the other player.} \\
\\ 
\texttt{PROPOSAL: \{'A08', 'A65', 'B96', 'A71', 'A25', 'B23', 'B19', 'C72', 'A94', 'A96', 'C79', 'B95'\}} \\
\\ 
\texttt{ARGUMENT: \{'This proposal uses exactly 8145 effort, the maximum allowed, ensuring we get full value from our shared limit. It includes 12 items that represent a balanced selection across different categories.'\}} \\
            }
        }
    }
    & \\ \\

    \theutterance \stepcounter{utterance}  
    & & & \multicolumn{4}{p{0.6\linewidth}}{
        \cellcolor[rgb]{0.9,0.9,0.9}{
            \makecell[{{p{\linewidth}}}]{
                \texttt{\tiny{[GM$\langle$P2]}}
                \texttt{STRATEGIC REASONING: \{'Let me analyze the opponent\textbackslash{}'s proposal. Their set: A08(963), A65(961), B96(901), A71(762), A25(738), B23(737), B19(730), C72(722), A94(704), A96(692), C79(731), B95(204). Total effort: 963+961+901+762+738+737+730+722+704+692+731+204 = 8145. This exactly hits the limit. For me, this gives importance: 1059+1057+997+858+834+833+826+818+800+788+827+300 = 9897. This is actually quite good for me {-} many of my high{-}importance items are included (A08, A65, B96 are my top 3, plus A71, A25, B23, B19, C72, A94, A96, C79 are all high value). However, I should check if I can do better. My top items by importance/effort ratio or absolute value: A75(32/128), C17(2/98), A28(44/140), A08(963/1059), A65(961/1057), B96(901/997), A71(762/858), A25(738/834), B23(737/833), C79(731/827), B19(730/826), C72(722/818), A94(704/800), A96(692/788), C10(637/733), A87(645/741), C82(633/729), B34(623/719), B37(510/606), B43(455/551), B16(405/501), C53(389/485), C51(369/465), A39(418/514), C84(307/403), A46(315/411), B15(250/346), B69(247/343), B64(230/326), B81(238/334), C83(216/312), B95(204/300), A97(156/252), C60(150/246), A88(175/271), A28(44/140), A75(32/128), C17(2/98). The opponent\textbackslash{}'s proposal excludes: C17, A28, A75, A88, A97, C60, C83, B81, B64, B69, B15, C84, A46, C51, B16, C53, B43, B37, B34, C82, A87, C10. Many of these are low importance for me. But C17(98), A28(140), A75(128) are very cheap. Could I make a better proposal? Let me try: keep A08, A65, B96, A71, A25, B23, B19, C72, A94, A96, C79 (same as theirs without B95), add C17, A28, A75. Effort: 8145{-}204+2+44+32 = 8019. That\textbackslash{}'s under limit. I could add more. Actually let me try to optimize: Start with must{-}haves: A08(963), A65(961), B96(901), A71(762), A25(738), B23(737), B19(730), C72(722), A94(704), A96(692), C79(731) = 7940. Remaining: 205. Can add: C17(2), A28(44), A75(32), A88(175) = 253, too much. C17+A28+A75 = 78, remaining 127. A97(156) too big. C60(150) too big. So C17+A28+A75+? Can\textbackslash{}'t add A88. Could add B95(204)? 78+204=282, too much. So best is C17+A28+A75 = 78, total 8018, remaining 127 unused. Or swap: remove something small, add B95? No B95 is only 300 importance. Actually let me try: remove B95 from their set (they have it), add C17, A28, A75. Effort change: {-}204+78 = {-}126, so 8019. Importance change for me: {-}300+98+140+128 = +66. So 9897+66 = 9963. Better! But I need to negotiate. Let me start by refusing and making a counter{-}proposal.'\}} \\
\\ 
\texttt{PROPOSAL: \{'A08', 'A65', 'B96', 'A71', 'A25', 'B23', 'B19', 'C72', 'A94', 'A96', 'C79', 'C17', 'A28', 'A75'\}} \\
\\ 
\texttt{ARGUMENT: \{'Your proposal uses the full limit but includes B95 which has relatively low value. By replacing B95 with C17, A28, and A75, we use 8019 effort (under limit) and gain three additional items. This gives us more total items with minimal effort waste while maintaining all the high{-}value selections.'\}} \\
\\ 
\texttt{REFUSE: \{'A08', 'A65', 'B96', 'A71', 'A25', 'B23', 'B19', 'C72', 'A94', 'A96', 'C79', 'B95'\}} \\
            }
        }
    }
     \\ \\

    \theutterance \stepcounter{utterance}  
    & & \multicolumn{4}{p{0.6\linewidth}}{
        \cellcolor[rgb]{0.9,0.9,0.9}{
            \makecell[{{p{\linewidth}}}]{
                \texttt{\tiny{[P1$\langle$GM]}}
                \texttt{PROPOSAL: \{'A08', 'A65', 'B96', 'A71', 'A25', 'B23', 'B19', 'C72', 'A94', 'A96', 'C79', 'C17', 'A28', 'A75'\}} \\
\\ 
\texttt{ARGUMENT: \{'Your proposal uses the full limit but includes B95 which has relatively low value. By replacing B95 with C17, A28, and A75, we use 8019 effort (under limit) and gain three additional items. This gives us more total items with minimal effort waste while maintaining all the high{-}value selections.'\}} \\
\\ 
\texttt{REFUSE: \{'A08', 'A65', 'B96', 'A71', 'A25', 'B23', 'B19', 'C72', 'A94', 'A96', 'C79', 'B95'\}} \\
            }
        }
    }
    & \\ \\

    \theutterance \stepcounter{utterance}  
    & \multicolumn{4}{p{0.6\linewidth}}{
        \cellcolor[rgb]{0.9,0.9,0.9}{
            \makecell[{{p{\linewidth}}}]{
                \texttt{\tiny{[P1$\rangle$GM]}}
                \texttt{STRATEGIC REASONING: \{'The opponent refused my proposal and made a counter{-}proposal. Let me analyze their proposal: A08(963)+A65(961)+B96(901)+A71(762)+A25(738)+B23(737)+B19(730)+C72(722)+A94(704)+A96(692)+C79(731)+C17(2)+A28(44)+A75(32) = 8019 effort. For me, this gives importance: 1059+1057+997+858+834+833+826+818+800+788+827+98+140+128 = 9037. My original proposal gave me 8997 importance. So their proposal actually gives me 40 MORE importance! They removed B95 (300 importance to me) and added C17(98)+A28(140)+A75(128) = 366 importance. This is better for me. However, I should check if I can negotiate for even more. But given that this already exceeds my original proposal's value and uses less effort (leaving 126 slack), I could try to add something. But wait {-} if I agree now, I get 9037. If I push back, they might discover I value those small items highly. Actually, let me check: can I add anything with 126 effort? Looking at items not in their proposal: B95(204) {-} too big, A39(418) {-} too big, etc. Nothing fits in 126. So this is actually optimal given their selection. I should agree {-} it benefits me more than my own proposal!'\}} \\
\\ 
\texttt{ARGUMENT: \{'Your proposal is well{-}reasoned. By swapping B95 for three smaller items, we gain more total items with less effort while preserving all the high{-}value selections. The 126 effort buffer also provides flexibility. This is an efficient use of our shared limit.'\}} \\
\\ 
\texttt{AGREE: \{'A08', 'A65', 'B96', 'A71', 'A25', 'B23', 'B19', 'C72', 'A94', 'A96', 'C79', 'C17', 'A28', 'A75'\}} \\
            }
        }
    }
    & & \\ \\

    \theutterance \stepcounter{utterance}  
    & & & \multicolumn{2}{p{0.3\linewidth}}{
        \cellcolor[rgb]{0.9,0.9,0.9}{
            \makecell[{{p{\linewidth}}}]{
                \texttt{\tiny{[GM$|$GM]}}
                \texttt{game successful} \\
            }
        }
    }
    & & \\ \\

    \theutterance \stepcounter{utterance}  
    & & & \multicolumn{2}{p{0.3\linewidth}}{
        \cellcolor[rgb]{0.9,0.9,0.9}{
            \makecell[{{p{\linewidth}}}]{
                \texttt{\tiny{[GM$|$GM]}}
                \texttt{end game} \\
            }
        }
    }
    & & \\ \\

\end{supertabular}
}

\end{document}
